Dall'estensione di $\mathbb{N}_0$ per includere gli elementi inversi rispetto alla somma,
si ottiene l'insieme dei numeri interi, denotato con $\mathbb{Z}$. Questo insieme gode delle stesse proprietà di
$\mathbb{N}_0$ rispetto alla somma, con l'aggiunta della proprietà dell'esistenza dell'elemento inverso.

\begin{property}
    \textbf{Esistenza dell'elemento inverso}
    \[\forall a \in \mathbb{Z}, \, \, \exists b \in \mathbb{Z} \mid a + b = b + a = 0.\]
\end{property}

\subsubsection{Gruppi abeliani}

La struttura algebrica $(\mathbb{Z}, +)$ risulta quindi un gruppo abeliano.
Una generica struttura algebrica $(\mathbb{G}, \cdot)$, formata da un insieme
$\mathbb{G} \neq \emptyset$ e un'operazione binaria "$\cdot$", è un gruppo se e solo se
soddisfa le seguenti proprietà.

\begin{property}
    \textbf{Chiusura di $\mathbb{G}$ su $\cdot$}
    \[\forall a, b \in \mathbb{G}, \; (a \cdot b) \in \mathbb{G}.\]
\end{property}
\begin{property}
    \textbf{Associativa}
    \[\forall a, b, c \in \mathbb{G}, \; (a \cdot b) \cdot c = a \cdot (b \cdot c).\]
\end{property}
\begin{property}
    \textbf{Esistenza dell'elemento neutro}
    \[\forall a \in \mathbb{G}, \; \exists e \in \mathbb{G} \mid a \cdot e = e \cdot a = a.\]
\end{property}
\begin{property}
    \textbf{Esistenza dell'elemento inverso}
    \[\forall a \in \mathbb{G}, \; \exists b \in \mathbb{G} \mid a \cdot b = b \cdot a = e.\]
\end{property}

\begin{observation}
    Non è necessaria la commutatività dell'operazione binaria per definire un gruppo
\end{observation}

Un gruppo si dice inoltre commutativo o abeliano se e solo se soddisfa la seguente proprietà.
\begin{property}
    \textbf{Commutativa}
    \[\forall a, b \in \mathbb{G}, \; a \cdot b = b \cdot a.\]
\end{property}

Affinché una generica struttura algebrica sia un gruppo abeliano,
è dunque necessario che soddisfi contemporaneamente tutte le proprietà sopra elenecate.

\subsubsection{Operazioni in $\mathbb{Z}$}

In $\mathbb{Z}$ è possibile definire l'operazione di prodotto "$\cdot$" nel seguente modo

\[
    \forall a, b \in \mathbb{Z}, \quad a \cdot b := \sum_{k=1}^{b} a.
\]

Si può dimostrare che l'operazione di prodotto gode delle seguenti proprietà.

\begin{property}
    \textbf{Commutativa}
    \[\forall a, b \in \mathbb{Z}, \;a \cdot b = b \cdot a.\]
\end{property}
\begin{property}
    \textbf{Associativa}
    \[\forall a, b, c \in \mathbb{Z}, \;(a \cdot b) \cdot c = a \cdot (b \cdot c).\]
\end{property}
\begin{property}
    \textbf{Esistenza dell'elemento neutro}
    \[\forall a \in \mathbb{Z}, \; a \cdot 1 = 1 \cdot a = a.\]
\end{property}
\begin{property}
    \textbf{Chiusura di $\mathbb{Z}$ rispetto al prodotto}
    \[\forall a, b \in \mathbb{Z}, \; (a \cdot b) \in \mathbb{Z}.\]
\end{property}

\begin{observation}
    All'interno di $\mathbb{Z}$ non può essere definito l'elemento inverso rispetto al prodotto.
\end{observation}
\begin{observation}
    L'operazione di prodotto può essere applicata anche all'insieme $\mathbb{N}_0$, nei confronti della quale
    anche quest'ultimo risulta chiuso.
\end{observation}